(1) 圆环磁通 \(\Phi_B=\pi a^2(B_0+bt)\)。沿 +z 观察取逆时针为环路正向，法拉第定律给
\[
\mathcal E=-\frac{d\Phi_B}{dt}=-\pi a^2b .
\]
故电流 \(I=\mathcal E/R=-\pi a^2b/R\)，沿 +z 观察为顺时针，其自生磁场反对正在增长的 +z 外场。

(2) 电阻热功率为
\[
P_J=I^2R=\frac{\pi^2a^4b^2}{R}>0.
\]
在 b 恒定且自感可忽略时电流恒定。若考虑有限自感，启停瞬态会出现额外储磁能项，本题已排除这一过程。

(3) 源装置维持外磁场变化时必须承担因感应电流而出现的附加功率；对电流环，电源对电荷的非保守电场做功率 \(\mathcal E I=\mathcal E^2/R=P_J\)，在静态环没有机械位移时全部耗散为热。故能量不由“磁力对电荷做功”直接产生，磁场变化依赖外部源的供能。

(4) 若 b<0，\(\mathcal E\) 和 I 同时反向，环产生的磁场阻碍外磁通减小；热功率 \(P_J\propto b^2\) 不变。这也验证楞次定律改变的是感应电流方向而非焦耳热正性。
